\documentclass[../NDCNNAI_main.tex]{subfiles}
\graphicspath{{\subfix{../fig/}}}

\begin{document}

\subsection{Универсальная аппроксимация в метрическом/банаховом пространстве}

\begin{definition}[Универсальная аппроксимация]
Пусть $(X, d)$ — метрическое пространство (часто банахово или нормированное). Класс функций $\mathcal{Y} \subset X$ обладает \textbf{свойством универсальной аппроксимации} (УА), если для любой функции $f \in X$ и любого $\varepsilon > 0$ существует функция $g \in \mathcal{Y}$ такая, что:
\[
d(f, g) < \varepsilon.
\]
В контексте искусственных нейронных сетей:
\begin{itemize}
    \item $X$ — пространство целевых функций (например, $C(K)$ или $L^p(\mu)$)
    \item $\mathcal{Y}$ — множество функций, реализуемых нейронной сетью с одним скрытым слоем:
    \[
    g(x) = \sum_{j=1}^{m} \alpha_j \sigma(w_j \cdot x + b_j),
    \]
    где $\sigma: \mathbb{R} \to \mathbb{R}$ — функция активации.
\end{itemize}
\end{definition}

\begin{remark}
Наличие УА означает, что сколь угодно сложную функцию можно приблизить с любой заранее заданной точностью с помощью достаточно большой нейронной сети определённой архитектуры.
\end{remark}

\subsection{Классические теоремы для $C(K)$}

\begin{theorem}[Цыбенко, 1989]
Пусть $K \subset \mathbb{R}^n$ — компактное множество, $C(K)$ — пространство непрерывных функций на $K$ с равномерной нормой $\|f\|_\infty = \sup_{x \in K} |f(x)|$. Пусть $\sigma: \mathbb{R} \to \mathbb{R}$ — непрерывная и \textbf{дискриминаторная} функция активации. Тогда множество функций вида
\[
G(x) = \sum_{j=1}^{N} \alpha_j \sigma(w_j \cdot x + \theta_j)
\]
плотно в $C(K)$.
\end{theorem}

\begin{definition}[Дискриминаторная функция]
Функция $\sigma$ называется \textbf{дискриминаторной} для $C(K)$, если для любой конечной знакопеременной регулярной борелевской меры $\mu$ на $K$ из условия
\[
\int_K \sigma(w \cdot x + b) \, d\mu(x) = 0 \quad \forall w \in \mathbb{R}^n, \; b \in \mathbb{R}
\]
следует, что $\mu \equiv 0$.
\end{definition}

\begin{proof}[Схема доказательства]
Предположим противное: линейная оболочка $S = \operatorname{span}\{\sigma(w\cdot x + b)\}$ не плотна в $C(K)$. Тогда:
\begin{enumerate}[label=\arabic*)]
    \item Существует собственное замкнутое подпространство $R = \overline{S} \subsetneq C(K)$.
    \item По теореме Хана–Банаха $\exists L \in C(K)^*$, $L \neq 0$, такой что $L\big|_R = 0$.
    \item По теореме Рисса–Маркова–Какутани $\exists \mu \in \mathcal{M}(K)$, $\mu \neq 0$:
    \[
    L(f) = \int_K f(x)\,d\mu(x) \quad \forall f \in C(K).
    \]
    \item В частности, для всех $w, b$:
    \[
    0 = L(\sigma(w\cdot x + b)) = \int_K \sigma(w\cdot x + b)\,d\mu(x).
    \]
    \item По дискриминаторности $\sigma$ имеем $\mu \equiv 0$, что противоречит $L \neq 0$.
\end{enumerate}
Следовательно, $S$ плотно в $C(K)$.
\end{proof}

\begin{example}
Дискриминаторными функциями являются:\\
\begin{itemize}
    \item Любая непрерывная сигмоидальная функция (логистическая сигмоида, гиперболический тангенс);
    \item Любая ограниченная измеримая неконстантная функция;
    \item ReLU (через комбинации можно построить ограниченные функции).
\end{itemize}
\end{example}
Одним из обобщений предыдущей теоремы является результат Хорника, Стинчкоба и Уайта.
\begin{theorem}[Хорник-Стинчкомб-Уайт, 1990]
Пусть $\sigma$ — непрерывная, ограниченная и неконстантная функция активации. Тогда для любого компакта $K \subset \mathbb{R}^n$ множество
\[
\mathcal{N}_\sigma = \left\{ \sum_{j=1}^{m} \alpha_j \sigma(w_j \cdot x + b_j) \right\}
\]
плотно в $C(K)$ в равномерной норме.
\end{theorem}

\begin{remark}
Эта теорема показывает, что:
\begin{itemize}
    \item УА достигается даже без требования сигмоидальности.
    \item Достаточно ограниченности и неконстантности активационной функции.
    \item Именно архитектура сети (один скрытый слой), а не конкретный вид активации, обеспечивает универсальность.
\end{itemize}
\end{remark}

\subsection{Понятие discriminatory/non-polynomial активации}
Ключевое значение в предыдущих теоремах играли следующие факты:

\begin{lemma}[Цыбенко]
Любая ограниченная, измеримая, сигмоидальная функция является дискриминаторной.
\end{lemma}

\begin{lemma}[Хорник]
\label{lem:hornik}
Любая ограниченная неконстантная функция является дискриминаторной (непрерывность не требуется).
\end{lemma}
Но это лишь достаточные условия, в 1993 Лешно, Лин, Пинкус и Шокен вывели критерий для функции активации, обеспечивающий УА для многослойного перцептрона с одним скрытым слоем в пространстве непрерывных функций.

\begin{theorem}[LLPS]
Пусть $\sigma$ — локально ограниченная и кусочно-непрерывная функция. Тогда
\[
\mathcal{N}_\sigma \text{ плотно в } C(\mathbb{R}^n) \text{ (равномерно на компактах)} \iff \sigma \text{ не полином п.в.} \]
\end{theorem}

\begin{proof}[Основные идеи]
\textbf{Необходимость:} если $\sigma$ — полином степени $d$, то $\mathcal{N}_\sigma$ состоит из полиномов степени $\leq d$, которые не плотны в $C(K)$.

\textbf{Достаточность:}
\begin{enumerate}[label=\arabic*)]
    \item Сведение к одномерному случаю через плотность хребтовых функций
    \item Рассмотрение свёрток $\sigma * \varphi$ с $\varphi \in C_c^\infty$
    \item Разбор двух случаев:
    \begin{itemize}
        \item Если $\exists \varphi: \sigma * \varphi$ не полином $\Rightarrow$ плотность через теорему Вейерштрасса
        \item Если $\forall \varphi: \sigma * \varphi$ --- полином $\Rightarrow$ $\sigma$ — полином почти всюду (противоречие)
    \end{itemize}
\end{enumerate}
\end{proof}

\begin{remark}
Теорема LLPS явно даёт необходимое и достаточное условие: активационная функция \textbf{не должна быть полиномом}. Это объясняет, почему ReLU, сигмоиды, тригонометрические функции работают, а квадратичная активация --- нет.
\end{remark}

\subsection{УА на $C(\mathbb{R}^n)$ с метрикой Фреше, связь с компактными множествами}
На $C(\mathbb{R}^n)$ равномерная норма неприменима (функции могут быть неограниченными). Используется \textbf{компактно-открытая топология}, задаваемая метрикой Фреше:

\begin{definition}[Метрика Фреше]
Для $f, g \in C(\mathbb{R}^n)$ определим:
\[
d(f, g) = \sum_{k=1}^\infty 2^{-k} \frac{\| (f - g) \cdot \mathbf{1}_{[-k,k]^n} \|_\infty}{1 + \| (f - g) \cdot \mathbf{1}_{[-k,k]^n} \|_\infty}.
\]
\end{definition}

\begin{remark}
Свойства метрики $d$:
\begin{itemize}
    \item $d$ --- метрика на $C(\mathbb{R}^n)$;
    \item $f_m \to f$ по $d$ $\iff$ $f_m \to f$ равномерно на каждом компакте;
    \item Веса $2^{-k}$ обеспечивают убывание влияния ошибок на больших кубах.
\end{itemize}
\end{remark}

\subsubsection*{УА на $C(\mathbb{R}^n)$}

\begin{theorem}[Пинкус-Стинчкомб, 1991]
Пусть $\sigma$ непрерывна и дискриминаторна. Тогда для любой $f \in C(\mathbb{R}^n)$ и $\varepsilon > 0$ существует нейронная сеть
\[
g(x) = \sum_{j=1}^m \alpha_j \sigma(w_j \cdot x + b_j)
\]
такая, что $d(f, g) < \varepsilon$.
\end{theorem}

\begin{proof}
    Приведём основные шаги доказательства теоремы.\\
    \textbf{Усечение хвоста:} Выбираем $M$ так, что $\sum_{k=M+1}^\infty 2^{-k} < \varepsilon/2$.\\
    \textbf{Аппроксимация на $K_M = [-M, M]^n$:} По теореме Цыбенко находим $g$ с
    \[
    \| (f - g) \cdot \mathbf{1}_{K_M} \|_\infty \leq \delta.
    \]
    \textbf{Оценка расстояния:}
    \[
    d(f, g) = \underbrace{\sum_{k=1}^M 2^{-k} \frac{\| (f - g)\mathbf{1}_{K_k} \|_\infty}{1 + \| \cdots \|_\infty}}_{\text{(I)}} + 
    \underbrace{\sum_{k=M+1}^\infty 2^{-k} \frac{\| \cdots \|_\infty}{1 + \| \cdots \|_\infty}}_{\text{(II)}},
    \]
    где (I) $\leq \frac{\delta}{1+\delta} \sum_{k=1}^M 2^{-k} \leq \frac{\delta}{1+\delta}$, а (II) $< \varepsilon/2$ по выбору $M$.
    
    Выбирая $\delta$ так, что $\frac{\delta}{1+\delta} < \varepsilon/2$ получаем итог: $d(f, g) < \varepsilon$.

\end{proof}

\begin{remark}
Ключевая идея: аппроксимация на всём $\mathbb{R}^n$ сводится к аппроксимации на последовательности компактов $K_k = [-k, k]^n$, исчерпывающих пространство. Метрика $d$ тут:
\begin{itemize}
    \item Взвешивает ошибки на компактах $K_k$ убывающими весами $2^{-k}$.
    \item Позволяет контролировать поведение на бесконечности.
    \item Согласована с компактно-открытой топологией.
\end{itemize}
\end{remark}
\end{document}